%% Cognitive Abstraction Layers (CAL) — Pre-Paper
%% arXiv submission draft — May 2026
%%
%% Compile sequence:
%%   pdflatex cal_prepaper.tex
%%   bibtex   cal_prepaper
%%   pdflatex cal_prepaper.tex
%%   pdflatex cal_prepaper.tex
%%
%% arXiv categories: cs.AI (primary), cs.LG, cs.IT
%%

\documentclass[11pt]{article}

% ---------- Encoding & fonts ----------
\usepackage[utf8]{inputenc}
\usepackage[T1]{fontenc}
\usepackage{lmodern}
\usepackage{microtype}

% ---------- Page geometry ----------
\usepackage[margin=1.1in, top=1.2in, bottom=1.2in]{geometry}

% ---------- Mathematics ----------
\usepackage{amsmath}
\usepackage{amssymb}
\usepackage{amsthm}

% ---------- Tables ----------
\usepackage{booktabs}
\usepackage{tabularx}
\usepackage{array}

% ---------- Lists ----------
\usepackage{enumitem}

% ---------- Code / verbatim ----------
\usepackage{listings}
\usepackage{xcolor}
\lstset{
  basicstyle=\small\ttfamily,
  breaklines=true,
  frame=single,
  rulecolor=\color{black!20},
  xleftmargin=1em,
  xrightmargin=0.5em,
}

% ---------- Bibliography ----------
\usepackage[numbers,sort&compress]{natbib}
\bibliographystyle{unsrtnat}

% ---------- Hyperlinks (last) ----------
\usepackage[colorlinks=true,
            linkcolor=black,
            citecolor=blue!60!black,
            urlcolor=blue!70!black]{hyperref}

% ---------- Theorem environments ----------
\theoremstyle{definition}
\newtheorem{hypothesis}{Hypothesis}
\newtheorem{definition}{Definition}
\theoremstyle{remark}
\newtheorem{remark}{Remark}

% ---------- Metadata ----------
\title{\textbf{Cognitive Abstraction Layers (CAL):\\
A Research Architecture for Hierarchical\\
Semantic Compression in AI Systems}}

\author{Juan Pablo Chancay\\
  \small Aural Syncro Research Lab\\
  \small \href{mailto:jpcpol@gmail.com}{jpcpol@gmail.com}}

\date{May 2026 \quad|\quad Version~1.2\\[4pt]
  \small Pre-print / Vision Paper --- arXiv submission candidate\\
  \small License: CC BY 4.0}

% ============================================================
\begin{document}
% ============================================================

\maketitle
\thispagestyle{empty}

% ---------- Abstract ----------
\begin{abstract}
\textbf{Thesis: scalable cognition requires hierarchical semantic
compression.} More precisely: inference scalability depends on
\emph{semantic state topology}, not raw context size. The quality
of AI governance scales with the structure of the compressed state
space --- with whether the compression preserves the topology of
governance-relevant states --- not with the number of tokens
processed. As AI systems generate outputs at scales that exceed
human supervisory capacity, the fundamental challenge of AI
governance shifts from generation to comprehension. We propose that
this challenge cannot be solved by increasing human review bandwidth
or improving artifact-level tooling, because both responses operate
at the wrong level of abstraction. We introduce the
\textbf{Cognitive Abstraction Layer (CAL) architecture}: a
five-level hierarchy spanning raw token streams~(L0) to autonomous
meta-inference on compressed semantic volumes~(L4). The core
invariant is \emph{semantic conservation}: each layer must compress
state while preserving sufficient causal and decisional structure
for effective reasoning at the layer above. We formalize this
requirement through \textbf{Semantic Information Density (SID)} ---
defined as the mutual information between the compressed
representation and the governance decision space, normalized by the
original. The central claim is that the quality of AI governance
scales with $\text{SID}_D$ at each layer boundary, not with context
window size. L1 is instantiated by existing large language model
inference but not formally structured as a governance layer; L2 is
formalized and empirically validated by the companion paper
TCO-L2~\cite{chancay2026tco}; L3 and L4 represent open research
directions with defined formal requirements and testable hypotheses.
\end{abstract}

\medskip
\noindent\textbf{Keywords:} cognitive governance, semantic compression,
tensor representation, hierarchical inference, mutual information,
Natural Cognitive Frontier, Semantic Information Density,
governance manifold

\tableofcontents
\newpage

% ============================================================
\section{Introduction}
% ============================================================

The dominant paradigm in AI-assisted software development positions
the human as a validator of AI-generated artifacts. As agent
pipelines mature, this paradigm is reaching a structural limit: the
volume of output generated by concurrent AI agents systematically
exceeds the human capacity for artifact-level review. Analysis of
code repositories with high proportions of AI-generated content
documents increasing defect rates despite larger output
volumes~\cite{gitclear2025}. The bottleneck is not generation --- it
is the human capacity to maintain systemic comprehension over the
outputs of multiple concurrent agents operating at different stages
of a shared pipeline.

Two common responses to this problem are: (1) reduce the human's
role further, replacing review with automated quality checks, and
(2) improve tooling around the existing paradigm --- better diffs,
richer IDEs, faster search. We argue that both responses are
incomplete, because they operate at the wrong abstraction level. The
problem is not that humans review too little, nor that their tools
are inadequate. The problem is that the information presented to
human supervisors is at the wrong abstraction level for the
cognitive operations they are being asked to perform.

This paper proposes a third response: a hierarchical architecture of
semantic abstraction layers that progressively compresses
AI-generated state into forms appropriate for increasingly autonomous
inference. We call this the \textbf{Cognitive Abstraction Layer
(CAL) architecture}. It has five levels, L0 through L4. The human
supervisor appears at Level~2 --- not as a reviewer of artifacts,
but as an orchestrator of system states. Below Level~2, generation
and local inference are automated; above Level~2, governance
inference is automated. At Level~2, the human operates at the
\textbf{Natural Cognitive Frontier (NCF)}~\cite{chancay2026tco}:
the abstraction level where cognitive demand is calibrated to human
capacity, decisions are expressible in natural language, and
judgment produces policy rather than correction.

The key question the architecture poses is not ``how much can humans
supervise?'' but ``what is the minimum semantic representation that
preserves the structure required for effective governance at each
level?'' This question is formalized through
\textbf{Semantic Information Density (SID)}: the quantity of
decision-relevant cognitive structure preserved per unit of
representation. The thesis of the CAL architecture is that
maximizing $\text{SID}_D$ at each layer boundary --- not expanding
context windows --- is the path to scalable AI governance.

The current paper is a research agenda. We do not report new
experimental results for L3 and L4. We provide: (a)~formal
definitions of each layer and the properties required at each
boundary; (b)~a summary of the L2 empirical anchor (TCO-L2);
(c)~concrete open problems with testable hypotheses for L3 and L4;
and (d)~the SID metric as a unifying measurement framework.

\subsection{Scope and Non-Goals}

CAL is a research architecture for AI pipeline governance in
multi-agent software development contexts. It is \emph{not}:

\begin{itemize}[leftmargin=1.5em]
  \item \textbf{A theory of general cognition or AGI.} The word
    ``cognitive'' in CAL refers to the supervisory operations that
    agents and humans perform over pipeline state, not to cognition
    in the cognitive science or AGI sense.

  \item \textbf{A proposal to eliminate human governance.} L2
    (TCO-L2) places the human orchestrator at the center of the
    architecture. L3 and L4 automate inference where human working
    memory is the bottleneck --- they do not eliminate human
    accountability.

  \item \textbf{A general inference efficiency claim.} The L4
    Efficiency Hypothesis applies specifically to governance
    inference over semantically compressed pipeline state. It does
    not claim that hierarchical abstraction improves efficiency for
    arbitrary tasks where semantic compression may destroy rather
    than preserve relevant structure.

  \item \textbf{A complete, deployed system.} L3 and L4 have
    defined formal requirements and open problems, not
    implementations.
\end{itemize}

The scope is: multi-agent AI pipelines with measurable quality
dimensions (as defined by $\varphi$ in Section~\ref{sec:l2}) and
known pipeline structure. Extension to other domains requires
domain-specific instantiation of $\varphi$ and fresh validation of
the L1$\to$L2 interface.

% ============================================================
\section{The CAL Architecture}
% ============================================================

The CAL architecture is a five-level hierarchy. Each level receives
compressed state from the level below, performs a defined
transformation, and outputs a further-compressed representation to
the level above. The architecture is bidirectional: state flows
upward (L0$\to$L4) and policy flows downward (L4$\to$L0). At~L2,
the human orchestrator is the bidirectional interface.

\begin{lstlisting}[caption={CAL architecture overview}]
L4  Meta-Inference Layer    M(V) -> {decisions, predictions}
     | M: V -> structured governance outputs
L3  Tensor Volume Layer     V(T_1, T_2, ..., T_n)
     | C: {T} -> V    [composition operator -- open problem]
L2  Cognitive Tensor Layer  T[d,i,j,k]       <- TCO-L2
     | f o phi: {A} -> T,  I: T -> {Omega, Delta, Rho, Xi}
     | Human orchestration at NCF (natural language policy)
L1  Semantic Feature Layer  {embeddings, relations, signals}
     | LLM inference: tokens -> features  [instantiated]
L0  Token / Artifact Space  raw AI-generated outputs
\end{lstlisting}

The core invariant at every boundary $L_k \to L_{k+1}$ is
\textbf{semantic conservation}: the compressed representation at
$L_{k+1}$ must preserve sufficient causal and decisional structure
for effective reasoning at $L_{k+1}$. A representation that fails
this invariant destroys governance-relevant information that cannot
be recovered at higher layers. Measuring the degree to which each
boundary satisfies this invariant is the role of SID
(Section~\ref{sec:sid}).

% ============================================================
\section{L0--L1: Foundation Layers}
% ============================================================

\subsection{L0: Token and Artifact Space}

L0 is the space of raw AI-generated outputs: source code,
configuration files, test results, documentation, deployment
manifests, system logs. L0 is the entry point to the abstraction
hierarchy and is not a contribution of CAL --- it is the existing
output of any AI-assisted development pipeline. Its key property is
\emph{volume asymmetry}: modern agentic pipelines generate artifacts
at rates that exceed human review capacity~\cite{gitclear2025},
motivating the entire hierarchy.

\subsection{L1: Semantic Feature Layer}

L1 is the space of semantic features extracted from L0 artifacts by
inference models: embeddings (dense vector representations of
artifact semantics), relational features (dependency graphs, call
graphs, architectural relationships), and quality signals (static
analysis outputs, test coverage, linting results).

\medskip
\noindent\textbf{Critical clarification --- L1 is instantiated, not
formalized.} Existing large language models generate L1-type features
as intermediate representations in service of their primary
generation and evaluation tasks. They do not expose a structured,
persistent L1 state that formally serves as input to a governance
layer. The CAL architecture does not propose building L1 from
scratch --- it formalizes the interface between L1 and L2, specifying
what properties L1 outputs must have to support effective L2
governance.

Formally, L1 produces a set of per-artifact feature representations
$\mathcal{F} = \{f_1, f_2, \ldots, f_n\}$ where each $f_i$ is a
local representation of a single artifact. L1 does not perform
cross-artifact, cross-temporal, or cross-agent aggregation --- these
are structurally the responsibility of L2.

% ============================================================
\section{L2: Cognitive Tensor Layer --- TCO-L2}
\label{sec:l2}
% ============================================================

L2 is formalized and empirically validated by the companion paper
TCO-L2~\cite{chancay2026tco}. This section summarizes its
contribution in the context of the CAL architecture.

\subsection{Vectorization: $\varphi$ --- L1 Features $\to$ Quality Vector}

\[
  \varphi: A \;\to\; \mathbf{V} \in [0,1]^{11}
\]

The vectorization function $\varphi$ maps each artifact to an
11-dimensional quality vector
$\mathbf{V} = (v_1, \ldots, v_{11})$, spanning: functional
correctness ($v_1$), architectural alignment ($v_2$), scalability
projection ($v_3$), security risk ($v_4$, inverted), observability
coverage ($v_5$), testability ($v_6$), maintainability ($v_7$),
technical debt ($v_8$, inverted), performance ($v_9$), evaluator
confidence ($v_{10}$), and anomaly score ($v_{11}$, inverted).

\medskip
\noindent\textbf{Epistemological note.} Dimensions $v_1, v_2, v_3,
v_9$ are \emph{supervisory estimators} --- heuristic semantic
signals generated by LLM assessment without universal ground truth.
Dimensions $v_4, v_6, v_7, v_8$ are \emph{verifiable ground truth}
--- outputs of deterministic static analysis (Bandit, Radon). This
distinction is critical for calibrating $\text{SID}_D$ at the
L1$\to$L2 boundary.

\subsection{Aggregation: $f$ --- Quality Vectors $\to$ Cognitive Tensor}

\[
  f: \{\mathbf{V}\} \;\to\; T[d,i,j,k] \in
  \mathbb{R}^{n \times |S| \times |A| \times |T_{\mathrm{idx}}|}
\]

The aggregation function $f$ builds a four-dimensional tensor indexed
by quality dimension~($d$), pipeline stage~($i$), agent~($j$), and
time index~($k$). The tensor's shared index structure
$[d,i,j,k]$ makes multidimensional supervisory operations
first-class.

Two experimental scenarios demonstrate the operational necessity of
this representation:

\begin{description}[leftmargin=0pt]
  \item[S3 --- Gradual technical debt.] Per-cycle degradation in
  $v_8$ of $-0.08$ falls within individual evaluator noise
  ($\sigma{=}0.07$, $\text{SNR}{=}1.14$). The 3-cycle cumulative
  delta ($-0.24$, $\text{SNR}{=}3.43$) reliably crosses the alert
  threshold. Monte Carlo simulation ($n{=}1000$) confirms:
  artifact-level detection rate 32.3\% vs.\ tensor $\Delta$
  detection rate 66.9\% ($+34.6$~pp gap). The tensor operation
  $T[v_8,:,:,0] - T[v_8,:,:,3]$ is one line; the SQL equivalent
  requires a self-join reconstructing the coordinate system.

  \item[S5 --- Inter-agent conflict.] Two agents produce opposing
  quality profiles ($v_4\,\Delta{=}0.65$, $v_6\,\Delta{=}0.50$).
  Per-artifact review requires the supervisor to hold both outputs in
  working memory simultaneously. The tensor $\varrho$ operation
  $|T[:,i,j_0,k] - T[:,i,j_1,k]|$ detects both conflicting
  dimensions automatically.
\end{description}

\subsection{Inference: $I$ --- Tensor $\to$ Governance Outputs}

\[
  I: T \;\to\; \{\Omega,\, \Delta,\, \varrho,\, \Xi\}
\]

\begin{center}
\begin{tabular}{@{}llll@{}}
  \toprule
  Symbol & Output & Threshold & Description \\
  \midrule
  $\Omega$ & Global state & stable $\geq 0.70$ &
    \texttt{stable / warning / critical} \\
  $\Delta$ & Trend & $|\text{slope}| > 0.05$ &
    Early-warning per dimension \\
  $\varrho$ & Systemic risk & $|\text{diff}| > 0.30$ &
    Inter-agent quality conflicts \\
  $\Xi$ & Recommendations & Top-3$\Delta$ + top-2$\varrho$ &
    Ranked by estimated $\partial\Omega/\partial\text{action}$ \\
  \bottomrule
\end{tabular}
\end{center}

\subsection{The Natural Cognitive Frontier (NCF)}

The NCF is the abstraction level at which cognitive demand is
calibrated to human capacity: decisions do not exceed working memory
bounds ($\leq 4$--$5$ chunks~\cite{cowan2001}), are expressible in
natural language, and produce policy rather than artifact correction.
At~L2, the human orchestrator reads $\{\Omega, \Delta, \varrho,
\Xi\}$ and injects policy $P_\mathrm{new}$ in natural language.

NCF operationalization in TCO-L2 uses four observable proxies:
NASA~Raw-TLX subscales (working memory saturation), correction log
consistency $\sigma$ across fault categories, accuracy variance
$\sigma(\alpha)$ across tasks T1$\to$T4, and interaction timer IQR
(attention fragmentation).

\subsection{L2 as the Empirical Anchor for Semantic Conservation}

The TCO-L2 experiment (between-subjects RCT, $n{=}40$ software
engineers) tests whether participants make more accurate governance
decisions from $T$ than from raw L0 artifacts. This is not only an
HCI claim --- it is the empirical validation that the L0$\to$L2
abstraction is semantically conservative.

\medskip
\noindent\textbf{Dual purpose of the experiment.} (1)~Validate that
the TCO-L2 dashboard reduces cognitive load and improves oversight
quality. (2)~Establish empirical evidence that tensor representation
preserves sufficient decision-relevant structure for effective
governance --- the first measured data point for $\text{SID}_D(L0 \to
L2)$. A positive result on H2 (Cohen's $d > 0.50$) is equivalent to
asserting $\text{SID}_D(L0 \to L2)$ above the baseline required for
governance.

% ============================================================
\section{L3: Tensor Volume Layer}
% ============================================================

L3 is the first layer above human working memory in the CAL
hierarchy. It operates on \textbf{tensor volumes}
$\mathbf{V}$ --- collections of L2 cognitive tensors aggregated
across sessions, pipelines, or time windows --- without requiring the
human as an inference substrate.

\subsection{Formal Definition}

Let $T^{(s)}$ denote the cognitive tensor for session~$s$. A tensor
volume is:
\[
  \mathbf{V} \;=\; C\!\left(\{T^{(s)}\}_{s=1}^{n}\right)
\]
where $C$ is the \textbf{composition operator} --- the central open
problem of L3. $C$ maps a collection of session tensors into a
unified higher-order structure preserving their causal and temporal
relationships.

\subsection{Requirements for the Composition Operator $C$}

For $\mathbf{V}$ to be semantically conservative with respect to L4
inference, $C$ must satisfy four properties:

\begin{enumerate}[label=\textbf{P\arabic*.},leftmargin=2.5em]
  \item \textbf{Causal preservation.} Causal relationships between
    quality dimensions recorded within any $T^{(s)}$ must be
    encodeable in $\mathbf{V}$. Using Pearl's do-calculus
    notation~\cite{scholkopf2021causal}, the interventional
    distribution $P(d_j \mid \mathrm{do}(d_i))$ must remain
    estimable from $\mathbf{V}$.

  \item \textbf{Temporal coherence.} The time indices of constituent
    tensors must be composable into a global temporal ordering.
    Cross-session drift must be representable in $\mathbf{V}$.

  \item \textbf{Dimensional stability.} The 11-dimensional quality
    vector structure must be preserved across the composition. L4
    must be able to apply the same inference operations
    $(\Omega, \Delta, \varrho, \Xi)$ to slices of $\mathbf{V}$
    that it applies to $T$.

  \item \textbf{Tractability.} $C$ must be computable for continuous
    pipelines. The size of $\mathbf{V}$ must not grow linearly with
    $n$ for large $n$.
\end{enumerate}

\subsection{Tractability and Mitigation Strategies}
\label{sec:l3-tractability}

The composition operator $C$ faces a combinatorial growth risk that
must be treated as a primary design constraint. A naive tensor
product $C = T^{(1)} \otimes T^{(2)} \otimes \cdots \otimes T^{(n)}$
grows exponentially in dimensionality, violating Property~P4 and
negating the L4 Efficiency Hypothesis. Four mitigation strategies
are candidates for the L3 research agenda:

\begin{enumerate}[leftmargin=2em]
  \item \textbf{Sparse tensor representation.} If agent-stage-di\-men\-sion
    combinations rarely co-occur, sparse formats (COO, CSR) reduce
    storage and computation to $\mathcal{O}(k)$ in the number of
    non-zero entries~$k$, where $k \ll n^{11}$.

  \item \textbf{Low-rank approximation.} Tucker and CP
    decompositions~\cite{kolda2009tensor} approximate high-dimensional
    tensors as products of smaller factor matrices. If the effective
    rank of $\mathbf{V}$ is bounded by the number of distinct
    governance-relevant patterns, low-rank $\mathbf{V}$ is tractable.

  \item \textbf{Selective retention (dynamic forgetting).} $C$ can
    implement a selective update rule --- analogous to gating in
    SSMs~\cite{gu2023mamba} --- discarding low-$\text{SID}_D$
    sessions and retaining high-$\text{SID}_D$ ones, keeping the
    effective~$n$ bounded.

  \item \textbf{Manifold projection.} If governance-relevant states
    lie on a low-dimensional manifold (see Section~\ref{sec:manifold}),
    $\mathbf{V}$ can be projected onto this manifold with bounded
    information loss.
\end{enumerate}

\subsection{Open Research Questions at L3}

\begin{itemize}[leftmargin=1.5em]
  \item Does $C$ belong to a known algebraic family (tensor product,
    contraction, CP/Tucker factorization)? What are the
    guarantees of each candidate?
  \item Does $\mathbf{V}$ exhibit stable \emph{attractor states}
    representing semantic invariants of the pipeline?
  \item Does $\mathbf{V}$ exhibit \emph{phase transitions} --- discontinuous
    changes in global governance state analogous to regime shifts in
    complex systems?
  \item How does the effective rank of $\mathbf{V}$ grow with $n$?
    Polynomial growth implies $\text{SID}_D$ decreases with pipeline
    scale --- a critical failure mode to characterize.
\end{itemize}

\subsection{Connections to Existing Research}

\textbf{State Space Models.} SSMs such as
Mamba~\cite{gu2023mamba} demonstrate that temporal sequences can be
compressed into fixed-size state representations with sub-quadratic
complexity while preserving long-range dependencies via selective
state spaces. The composition operator $C$ in L3 is a semantic
analog of the SSM state transition, operating on semantically
structured quality tensors rather than raw token sequences.

\medskip
\noindent\textbf{Renormalization Group Theory.} In statistical
physics, renormalization reduces the degrees of freedom in a system
while preserving macroscopic behavior~\cite{wilson1971renormalization}.
L3 is a semantic coarse-graining: it reduces the dimensionality of
the cognitive state space while preserving macroscopic governance
properties (trends, conflicts, anomalies) that matter for L4
inference.

\medskip
\noindent\textbf{Tensor Decomposition.} Tucker, CP, and tensor
train/ring formats~\cite{kolda2009tensor} offer tractable low-rank
approximations and are candidates for $C$ when exact composition is
intractable.

\medskip
\noindent\textbf{Causal Representation Learning.} The do-calculus
requirement for $C$ (Property P1) connects to causal representation
learning~\cite{scholkopf2021causal}. Two methods are directly
applicable: (a)~\emph{causal graph priors}, encoding known causal
relationships between quality dimensions as structural constraints on
$C$; and (b)~\emph{causal regularization}, penalizing representations
where interventional distributions $P(d_j \mid \mathrm{do}(d_i))$ are
not recoverable from $\mathbf{V}$. Causal graph priors derived from
domain knowledge about software pipeline dynamics are the more
tractable near-term approach, since causal discovery in compressed
spaces is computationally expensive~\cite{scholkopf2021causal}.

\subsection{The Governance Manifold Hypothesis}
\label{sec:manifold}

The tractability mitigations in Section~\ref{sec:l3-tractability} are
most effective when the following foundational hypothesis holds.

\begin{hypothesis}[Governance Manifold]
The set of governance-relevant pipeline states is a low-dimensional
manifold $\mathcal{M}_{\mathrm{gov}}$ embedded in the full tensor
space
$\mathbb{R}^{n \times |S| \times |A| \times |T_{\mathrm{idx}}|}$.
The intrinsic dimensionality $\dim(\mathcal{M}_{\mathrm{gov}})$ is
bounded by the number of distinct governance-relevant failure
patterns the pipeline can produce, not by the full tensor
dimensionality.
\end{hypothesis}

If true, the consequences propagate through the entire architecture:
\begin{itemize}[leftmargin=1.5em]
  \item \textbf{Tractability.} $\mathbf{V}$ needs only to represent
    $\mathcal{M}_{\mathrm{gov}}$, not the ambient space.
    $\dim(\mathcal{M}_{\mathrm{gov}})$ may be orders of magnitude
    smaller than the ambient dimension.

  \item \textbf{SID geometry.} $\text{SID}_D(L_k \to L_{k+1}) > \theta$
    becomes a topological question: does $C$ preserve the connectivity
    and local curvature of $\mathcal{M}_{\mathrm{gov}}$?

  \item \textbf{L4 Efficiency.} If $M(\mathbf{V})$ operates on
    $\mathcal{M}_{\mathrm{gov}}$, then
    $\text{Cost}(M(\mathbf{V})) = \mathcal{O}(\dim(\mathcal{M}_{\mathrm{gov}}))$,
    decoupled from $n$. The L4 Efficiency Hypothesis becomes a
    consequence, not an independent claim.

  \item \textbf{Semantic collapse framing.} Collapse
    (Section~\ref{sec:sid-openprob}) corresponds precisely to $C$
    reducing $\mathcal{M}_{\mathrm{gov}}$ below its intrinsic
    dimensionality.
\end{itemize}

\medskip
\noindent\textbf{Testability at L2.} The hypothesis is testable on
the TCO-L2 corpus without new data collection. If quality vector
trajectories across scenarios S1--S5 can be embedded in a
low-dimensional space via UMAP~\cite{mcinnes2018umap} or Isomap
while preserving inter-scenario distances, this supports a
low-dimensional governance manifold at~L2. The embedding
dimensionality is the first measurable proxy for
$\dim(\mathcal{M}_{\mathrm{gov}})$.

This positions topology --- not compression ratio --- as the central
mathematical object of CAL: the right question is not ``by how much
did we compress?'' but ``did we preserve the topology of the
governance manifold?''

% ============================================================
\section{L4: Meta-Inference Layer}
% ============================================================

L4 performs inference on tensor volumes $\mathbf{V}$ rather than on
individual tensors $T$ or artifacts $A$. It is the fully automated
inference layer: no human working memory is required as the inference
substrate.

\subsection{Formal Definition}

\[
  M: \mathbf{V} \;\to\; \{\text{decisions},\;
  \text{predictions},\; \text{adaptations}\}
\]

Because $\mathbf{V}$ (by Properties P1--P4 of $C$) encodes causal
relationships, temporal coherence, and dimensional stability, $M$
can perform reasoning not possible at lower layers:

\begin{itemize}[leftmargin=1.5em]
  \item \textbf{Causal discovery:} identifying which quality
    dimensions causally influence others, across sessions.
  \item \textbf{Drift prediction:} predicting future quality
    trajectories from historical volume dynamics, before the signal
    crosses the L2 alert threshold.
  \item \textbf{Cross-pipeline conflict detection:} identifying
    conflicts between architectural strategies spanning multiple
    sessions.
  \item \textbf{Policy generation:} producing governance policies
    from historical patterns for human review at L2 or direct
    injection in high-confidence, low-stakes cases.
\end{itemize}

\subsection{The L4 Efficiency Hypothesis}

\begin{hypothesis}[L4 Efficiency]
There exists an inference architecture such that the cost of
$M(\mathbf{V})$ scales with $\kappa(\mathbf{V})$ --- the structural
complexity of $\mathbf{V}$ (e.g., its effective rank, the entropy
of its attractor distribution, the size of its causal graph) ---
where $\kappa(\mathbf{V})$ grows significantly more slowly than
$\mathcal{O}(n^2)$ in $n$ (the raw artifact count of the underlying
L0 space):
\[
  \text{Cost}(M(\mathbf{V})) = \mathcal{O}(\kappa(\mathbf{V}))
  \quad\text{where}\quad
  \kappa(\mathbf{V}) \ll \mathcal{O}(n^2).
\]
\end{hypothesis}

This is not a claim that meta-inference is exponentially cheaper than
token-level attention in all cases. It is a claim that if semantic
abstraction is semantically conservative, inference cost can be
decoupled from the volume of the original artifact stream and coupled
instead to the structural complexity of the compressed governance
state. This connects to the observation that current transformer
architectures re-evaluate redundant relationships at every forward
pass~\cite{tay2022efficient}; a hierarchy of semantic abstractions
would eliminate this redundancy before inference.

Empirical testing requires: (a)~a concrete definition of $C$ and
$\kappa(\mathbf{V})$; (b)~a comparison of $M(\mathbf{V})$ cost vs.\
equivalent flat-context inference on raw artifacts; and
(c)~measurement of governance accuracy under both approaches.
TCO-L2 provides the L2 baseline for~(c).

\subsection{Open Research Questions at L4}

\begin{itemize}[leftmargin=1.5em]
  \item What is the minimal $\mathbf{V}$ that supports $M$ for a
    given class of governance decisions?
  \item Can $M$ be trained end-to-end from labeled governance
    outcomes, or does it require explicit causal structure in
    $\mathbf{V}$ as supervision?
  \item What are the failure modes of $M$ when $C$ fails to preserve
    causal structure? How does governance accuracy degrade?
  \item At what scale (agents, sessions, dimensions) does
    $M(\mathbf{V})$ outperform equivalent flat-context LLM inference
    on governance tasks?
\end{itemize}

% ============================================================
\section{Semantic Information Density}
\label{sec:sid}
% ============================================================

The CAL architecture requires a metric applicable at every layer
boundary to assess whether the abstraction is semantically
conservative. Without such a metric, the claim that each layer
``preserves decision-relevant structure'' is unfalsifiable. We
propose \textbf{Semantic Information Density (SID)}.

\subsection{Formal Definition}

Let $R_k$ denote the representation at layer $L_k$, and let $D$ be
the set of governance decisions relevant to the pipeline. Define the
\textbf{decision-relevant information content} of $R_k$ as the
mutual information between $R_k$ and $D$:
\[
  I_\mathrm{decision}(R_k)
    \;=\; I(D\,;\, R_k)
    \;=\; H(D) - H(D \mid R_k).
\]

\begin{definition}[Semantic Information Density]
\[
  \mathrm{SID}_D(L_k \to L_{k+1})
    \;=\; \frac{I(D\,;\, R_{k+1})}{I(D\,;\, R_k)}
    \;\in\; [0,\,1].
\]
\end{definition}

A $\text{SID}_D$ of $1.0$ means $R_{k+1}$ is as informative as
$R_k$ for governance decisions. A $\text{SID}_D$ below threshold
$\theta$ means the compression has destroyed governance-relevant
information.

\medskip
\noindent\textbf{Mathematical guarantee: $\text{SID}_D \in [0,1]$.}
By the data processing inequality~\cite{cover2006elements}, for the
Markov chain $D \to R_k \to R_{k+1}$ (where $R_{k+1}$ is produced
by applying $C$ to $R_k$):
\[
  I(D\,;\, R_{k+1}) \;\leq\; I(D\,;\, R_k),
\]
guaranteeing $\text{SID}_D \leq 1$ for any valid compression,
without additional assumptions. Information destroyed at a layer
boundary cannot be recovered at higher layers --- making
$\text{SID}_D$ a strict, irreversible measure of information loss.

\medskip
\noindent\textbf{Causal structure.} A critical dimension of SID not
captured by mutual information alone is causal structure
preservation. A compression may preserve high $\text{SID}_D$ for
contemporaneous decisions while destroying the interventional
distributions $P(d_j \mid \mathrm{do}(d_i))$ needed for drift
prediction and temporal governance --- a silent failure mode.

\medskip
\noindent\textbf{SID is contextual, not absolute.} Because
$I(D;\,R_k)$ depends on the governance decision set $D$, SID is
not a single scalar property --- it is a family indexed by governance
class:
\[
  \mathrm{SID}_D(L_k \to L_{k+1})
    \quad \text{for a specific governance class } D.
\]
All SID claims in this paper are relative to the software
development governance class defined by the TCO-L2 experimental
scenarios~(S1--S5). Extension to other domains requires fresh
operationalization of~$D$.

\subsection{Operationalization at L2}

The TCO-L2 experiment provides the first operational measurement of
$\text{SID}_D$. Let $\alpha_2$ be the governance accuracy of the
experimental group (operating at L2) and $\alpha_0$ the accuracy of
the control group (operating at L0). A bounded estimate:
\[
  \widehat{\mathrm{SID}}_D(L0 \to L2)
    \;\approx\; \frac{\alpha_2}{\alpha_\mathrm{max}},
\]
where $\alpha_\mathrm{max}$ is the accuracy achievable with perfect
information. H2 (Cohen's $d > 0.50$) is equivalent to asserting
$\text{SID}_D(L0 \to L2)$ substantially above baseline.

\subsection{The Open Problem: SID Without a Human Reference}
\label{sec:sid-openprob}

At L2, the human provides the calibration signal. At L3 and L4,
this signal is unavailable. Three candidate approaches:

\begin{enumerate}[leftmargin=2em]
  \item \textbf{Downstream outcome tracking.} Measure whether
    $M(\mathbf{V})$ decisions lead to measurable pipeline
    improvements over time.

  \item \textbf{Synthetic benchmarks with known causal structure.}
    Generate pipelines with known ground-truth causal graphs; test
    whether $M(\mathbf{V})$ recovers the correct relationships.

  \item \textbf{Comparison against L2 human decisions on held-out
    cases.} For novel governance scenarios, compare $M(\mathbf{V})$
    against expert human decisions at L2 as a delayed reference.
\end{enumerate}

\medskip
\noindent\textbf{Semantic collapse as a named failure mode.}
\textbf{Semantic collapse} occurs when $C$ preserves acceptable
$\text{SID}_D$ for each individual governance class~$D$ while
simultaneously destroying the diversity of the governance state
space --- collapsing distinctions between states relevant to
\emph{different} decision classes. This is analogous to mode
collapse in generative models: local fidelity is maintained, global
coverage is lost. Semantic collapse is particularly dangerous
because it passes contemporaneous SID checks while producing
systematic failures on novel or rare governance scenarios. Three
mitigations: (a)~\emph{entropy preservation constraints}, requiring
$H(\mathbf{V})$ to remain within a specified band of $H(T)$;
(b)~\emph{attractor separation metrics}, requiring that distinct
attractor states remain metrically separated in $\mathbf{V}$;
(c)~\emph{governance diversity benchmarks}, test suites deliberately
including rare failure modes to detect collapsed representations.

The Governance Manifold Hypothesis provides the geometric
interpretation: collapse corresponds to $C$ reducing
$\dim(\mathcal{M}_{\mathrm{gov}})$ below the intrinsic
dimensionality of the governance manifold.

% ============================================================
\section{Research Agenda}
% ============================================================

\begin{table}[h]
\centering
\caption{CAL research agenda by layer}
\label{tab:agenda}
\renewcommand{\arraystretch}{1.2}
\begin{tabularx}{\textwidth}{@{}lllX@{}}
\toprule
\textbf{Layer} & \textbf{Status} & \textbf{Target} &
  \textbf{Key open problem / testable hypothesis} \\
\midrule
L0 & Production & --- &
  Baseline; no open problem \\
L1 & Instantiated & --- &
  Formal L1$\to$L2 interface; LLM signals have $\text{SID}_D > \theta$ \\
\textbf{L2} & \textbf{TCO-L2 (CHI~2027)} & \textbf{CHI~2027} &
  NCF operationalization; H2: Cohen's $d > 0.50$;
  $\text{SID}_D(L0{\to}L2) > 0.70$ \\
L3 & Open & NeurIPS/ICML &
  Composition operator $C$ (P1--P4);
  $\text{SID}_D(L2{\to}L3) > \theta$ \\
L4 & Open & NeurIPS/ICML &
  $M$ architecture; $\text{Cost}(M(\mathbf{V})) = \mathcal{O}(\kappa(\mathbf{V}))$
  where $\kappa(\mathbf{V}) \ll \mathcal{O}(n^2)$ \\
SID & Partial (L2) & Cross-layer &
  Measure $\text{SID}_D$ without human reference; define $D$ for
  non-software domains \\
\bottomrule
\end{tabularx}
\end{table}

% ============================================================
\section{Positioning Within Existing Research}
% ============================================================

\textbf{AI Observability / MLOps.} Datadog, New Relic, and
OpenTelemetry monitor operational health (latency, errors,
resources). TCO-L2 monitors semantic health (quality, debt,
architectural alignment, inter-agent conflicts). CAL extends this to
hierarchical governance.

\medskip
\noindent\textbf{Hierarchical Reinforcement Learning.} HRL
decomposes long-horizon tasks across temporal abstractions (options
framework). CAL applies analogous hierarchical abstraction to
governance, not to task execution.

\medskip
\noindent\textbf{State Space Models.} SSMs~\cite{gu2023mamba}
compress temporal sequences with sub-quadratic complexity. L3
requires a semantic analog of the SSM state transition for
structured quality tensors.

\medskip
\noindent\textbf{Global Workspace Theory.} GWT~\cite{baars1988cognitive}
proposes a shared workspace integrating specialized cognitive
processes. The L2 tensor is a computational analog: a shared index
structure integrating quality signals from multiple specialized
evaluation agents.

\medskip
\noindent\textbf{Situation Awareness.} Endsley's SA
model~\cite{endsley1995} (perception $\to$ comprehension $\to$
projection) maps to L0$\to$L1$\to$L2 in the CAL hierarchy. TCO-L2
is an operationalization of SA Levels~2--3 for AI pipeline
supervision.

% ============================================================
\section{Conclusion}
% ============================================================

The CAL architecture proposes a reframing of the AI governance
problem: not as a question of how much humans can supervise, but as
a question of what abstraction level at each layer boundary
maximizes $\text{SID}_D$ for governance decisions. TCO-L2 is the
empirical anchor --- it validates, via controlled experiment, that
the L0$\to$L2 abstraction is semantically conservative enough for
effective human orchestration at the Natural Cognitive Frontier.
That validation is the theoretical foundation for L3 and L4: if a
human can govern effectively from a tensor representation, an
automated inference layer operating on tensor volumes can do so
without the human working memory constraint.

The central open problems --- the composition operator $C$ for L3,
the measurement of $\text{SID}_D$ without human reference, the L4
Efficiency Hypothesis, and the Governance Manifold Hypothesis ---
define a coherent research agenda with concrete testable hypotheses
at each step. The claim is not that these problems are solved, but
that they are the right problems to solve: progress on each one
directly advances scalable AI governance.

% ============================================================
\bibliography{cal_prepaper}
% ============================================================

\end{document}
